\section{Bewertung der Effektivität}
\label{ch:Diskussion:Effektivität}

In \autoref{ch:Ergebnisse:sec:Hypothesen_Auswertung:sub:Effektivität} wurde aufgezeigt, dass die Fehlerquote des Dash\-boards im Schnitt höher und somit schlechter ist als die des alten Systems, dies jedoch keinen signifikanten Unterschied darstellt (s. \autoref{fig:usability_test:fehlerquote}).

Eine Ungenauigkeit in den Aufgaben 3.1, 3.2, 4.1 und 4.2 hat dazu geführt, dass jeweils zwei Antworten im alten System akzeptiert werden mussten. Während das Ziel auf dem Dash\-board eindeutig identifizierbar abgelesen werden kann, ist in Aufgabe 4.1 und 4.2 der Begriff 'noch nicht gelaufen' im alten System verwechselt worden mit der Fehleranzeige, dass der Test 'schief gelaufen ist'. In den Aufgaben 3.1 und 3.2 wurde eine Testkategorie gemeint zum Zusammenzählen, jedoch hat einer der Testfälle den gleichen Namen. Beide Fälle haben jedoch keine Auswirkung auf das Ergebnis.

Ein Grund, wie die Quote entstanden ist, welcher auch von den Testpersonen genannt wurde, ist, dass das Dash\-board anders aufgebaut ist, als das alte System. Wie in \autoref{ch:Dashboard:sec:Aufbau} erwähnt, wurden die Zeilen und Spalten der Kategorien und Inhalte getauscht, was zu Verwirrung beim Finden der richtigen Daten führte.

Dass das alte System in vier von sechs Fragen keine Fehler erzeugt hat, lässt sich damit erklären, dass die meisten Nutzer dieses System seit über 30 Jahren nutzen. Zum Lernerfolg wird die Wiederholung einer Tätigkeit benötigt (s. \autoref{ch:SdF:UI:sec:menschenzentriertes_Design}). Dabei kann bei dem neuen System keine Vorerfahrung verzeichnet werden. Das System war für alle Beteiligten neu, zumal sich nur eine Person die neue Umgebung angesehen hat, bevor die Aufgaben begonnen wurden (s. \autoref{ch:Validierung:sec:Stichprobe}).

Der Anstieg der Fehlerquote im alten System bei Aufgabe 4 zeigt auf, dass dort Werte für das Ergebnis zusammengezählt werden mussten, während die geforderte Antwort vom Dash\-board lediglich abzulesen  waren. Es wurde sich häufiger verzählt oder verrechnet. Zusätzlich wurden Aussagen getroffen, dass einzelne Personen nicht mit dem System im Konkreten arbeiten, wodurch das Wissen darüber eingeschränkt ist.

Dies trägt auch zur nachträglichen Einschätzung bei, dass die Aufgaben 3 und 4 wenig Relevanz für den Alltag haben. Die Aufgaben 3.1 und 3.2 befassen sich mit der Release-Version, welche kein Teil des Alltagsgeschäft ist. Die Aufgaben 4.1 und 4.2 handeln von einer Zählaufgabe im alten System, wobei im Alltag eher wichtig ist, ob es noch Tests gibt, die nicht gelaufen sind und nicht wie viele davon. Daher werden diese Aufgaben aus der Auswertung rausgenommen.

Mit den objektiven Ergebnissen der Aufgaben 1 und 2 kann die Hypothese $H_{11}$ nun doch bestätigt werden.

In Bezug auf die Auswertung der Nutzerbefragung mit den Fragen zu den Prinzipien Aufgabenangemessenheit und Steuerbarkeit, sind in den Daten jeweils zwei Ausreißer aufgefallen (s. \autoref{tab:Nutzerbefragung:rohe_Daten}).
Bei den Fragen 2 und 3 haben zwei Personen entgegen allen weiteren Teilnehmern mit verdrehten Werten geantwortet. Sie sind die einzigen beiden, die das Dash\-board negativ und das alte System positiv bewerten. Die beiden Testpersonen fallen zusätzlich auch in anderen Bereichen noch auf.

Da der Verdacht auf einen Fehler zweiter Art, dem Beibehalten der Nullhypothese trotz Widersprüchen, besteht, wurden aus diesem Grund die beiden Signifikanztests noch einmal ohne diese Werte vorgenommen (s. \autoref{Add:Pruemper:neuer_ttest:frage_2_und_3}). Dadurch ist auch bei diesen Fragen ein signifikanter Unterschied zugunsten des Dash\-boards festzustellen.

Die Hypothese $H_{11}$ kann somit in Kombination mit den anderen Fragen der beiden Prinzipien auch auf Grundlage der subjektiven Ergebnisse bestätigt werden.

Diese Auffälligkeit ist auch in weiteren Datensätzen teilweise mit extremeren Differenzen zu beobachten (s. \autoref{fig:hist_prümper_all}), hat allerdings den p-Wert unter Ausschluss der Datensätze lediglich verringert und liegt weiterhin unterhalb des Signifikanzniveaus.

Zusammenfassend scheinen die Testpersonen das Gefühl zu haben, dass sie das Dash\-board gut nutzen können. Die eigenen Fehler oder Slips (s. \autoref{ch:SdF:UI:sec:menschenzentriertes_Design}) sind ihnen dabei nicht aufgefallen. Die Hypothesenbewertung auf Grundlage der subjektiven und objektiven Daten bilden ein schlüssiges Bild.
